\section{从波长关系到束缚态方程}
\subsection{经典轨迹和量子态的不同描述}
玻尔模型把经典圆轨道与量子化条件相结合，能解释氢样原子的主要能量尺度，却没有给出一般波函数和测量概率。处理波函数节点、势垒穿透和一般角动量态，需要量子态的描述。更完整的一维描述使用复波函数$\psi(x,t)$，$|\psi|^2$为位置概率密度，满足归一化$\int|\psi|^2\dd x=1$。

自由粒子的德布罗意关系$p=\hbar k$、$E=\hbar\omega$与非相对论色散$E=p^2/(2m)$相容于方程$i\hbar\partial_t\psi=-\hbar^2\partial_x^2\psi/(2m)$。加入给定势能$V(x)$，非相对论量子理论采用
\[
 i\hbar\partial_t\psi=\left[-\frac{\hbar^2}{2m}\partial_x^2+V(x)\right]\psi.
\]
这是一条量子动力学的基本方程；前面的自由粒子关系只说明它与德布罗意关系相容。对不含时势，取$\psi=\varphi(x)\ee^{-iEt/\hbar}$，则$\partial_t\psi=(-iE/\hbar)\psi$、$\partial_x^2\psi=\varphi''\ee^{-iEt/\hbar}$。代入并约去共同的时间因子，得到$-\hbar^2\varphi''/(2m)+V\varphi=E\varphi$。允许能量还受边界条件约束，不能只由局部$p^2=2m(E-V)$确定。

\subsection{边界连续性与允许区}
有限势跳变处，积分定态方程跨过一个无穷小区间。只要势没有狄拉克型奇点，积分中的有限项趋零，得到导数连续；波函数本身也连续，避免二阶导数产生未被势抵消的奇异项。无限高硬壁是极限边界，波函数在壁外为零、壁处为零，壁内导数不需要与外侧零函数继续相等。

在$E>V$区，局部波数实，波函数可振荡；$E<V$区，局部波数纯虚，波函数可以指数衰减或增长。有限势垒内的非零尾部不表示经典粒子拥有负动能，而是波函数满足另一类局部微分方程。对半无限禁区排除增长解，对有限禁区则一般保留两支以匹配两端。

\Diagram{fig:m4allowed}{线性势的艾里解平滑连接允许区与禁区}{\draw[axis](-.4,0)--(6,0)node[right]{$x$};\draw[axis](0,-1.2)--(0,1.9)node[above]{$\varphi$};\draw[guide](3,-1.2)--(3,1.8);\draw[curve]plot coordinates {(0.00000,-0.9470357) (0.05046,-0.9038078) (0.10092,-0.8537957) (0.15138,-0.7974842) (0.20183,-0.7353909) (0.25229,-0.6680602) (0.30275,-0.5960570) (0.35321,-0.5199612) (0.40367,-0.4403617) (0.45413,-0.3578511) (0.50459,-0.2730203) (0.55505,-0.1864540) (0.60550,-0.0987258) (0.65596,-0.0103941) (0.70642,0.0780016) (0.75688,0.1659438) (0.80734,0.2529406) (0.85780,0.3385279) (0.90826,0.4222713) (0.95872,0.5037687) (1.00917,0.5826506) (1.05963,0.6585822) (1.11009,0.7312630) (1.16055,0.8004281) (1.21101,0.8658473) (1.26147,0.9273255) (1.31193,0.9847018) (1.36239,1.0378485) (1.41284,1.0866705) (1.46330,1.1311037) (1.51376,1.1711138) (1.56422,1.2066944) (1.61468,1.2378660) (1.66514,1.2646733) (1.71560,1.2871841) (1.76606,1.3054872) (1.81651,1.3196902) (1.86697,1.3299179) (1.91743,1.3363102) (1.96789,1.3390201) (2.01835,1.3382122) (2.06881,1.3340601) (2.11927,1.3267457) (2.16972,1.3164565) (2.22018,1.3033846) (2.27064,1.2877250) (2.32110,1.2696741) (2.37156,1.2494285) (2.42202,1.2271835) (2.47248,1.2031323) (2.52294,1.1774648) (2.57339,1.1503667) (2.62385,1.1220186) (2.67431,1.0925954) (2.72477,1.0622656) (2.77523,1.0311908) (2.82569,0.9995253) (2.87615,0.9674155) (2.92661,0.9349999) (2.97706,0.9024089) (3.02752,0.8697645) (3.07798,0.8371804) (3.12844,0.8047618) (3.17890,0.7726057) (3.22936,0.7408008) (3.27982,0.7094279) (3.33028,0.6785597) (3.38073,0.6482616) (3.43119,0.6185912) (3.48165,0.5895995) (3.53211,0.5613302) (3.58257,0.5338209) (3.63303,0.5071029) (3.68349,0.4812018) (3.73394,0.4561376) (3.78440,0.4319255) (3.83486,0.4085756) (3.88532,0.3860938) (3.93578,0.3644820) (3.98624,0.3437383) (4.03670,0.3238576) (4.08716,0.3048314) (4.13761,0.2866488) (4.18807,0.2692963) (4.23853,0.2527582) (4.28899,0.2370171) (4.33945,0.2220536) (4.38991,0.2078474) (4.44037,0.1943766) (4.49083,0.1816186) (4.54128,0.1695500) (4.59174,0.1581466) (4.64220,0.1473842) (4.69266,0.1372379) (4.74312,0.1276832) (4.79358,0.1186951) (4.84404,0.1102491) (4.89450,0.1023207) (4.94495,0.0948859) (4.99541,0.0879210) (5.04587,0.0814028) (5.09633,0.0753088) (5.14679,0.0696167) (5.19725,0.0643052) (5.24771,0.0593535) (5.29817,0.0547415) (5.34862,0.0504499) (5.39908,0.0464601) (5.44954,0.0427541) (5.50000,0.0393148)};\node[above]at(1.2,1.8){允许区};\node[above]at(4.5,1.8){禁区};}
图\ref{fig:m4allowed}画出线性势中的同一个艾里函数解：转点左侧振荡，右侧衰减，函数及其导数在转点处连续。两区振幅由这个连接条件确定，不能分别任取。有限方阱解中的$\arctan(k/\kappa)$边界相位同样来自函数和导数的同时匹配。

\section{WKB相位、振幅和量子偏移}
令定态波函数$\varphi=A\ee^{iS/\hbar}$。先求一次导数$\varphi'=(A'+iAS'/\hbar)\ee^{iS/\hbar}$，再对括号与指数分别求导，得到
\[
 \varphi''=\left[A''+\frac{2i}{\hbar}A'S'
 +\frac{i}{\hbar}AS''-\frac1{\hbar^2}A(S')^2\right]\ee^{iS/\hbar}.
\]
当$A,S$相对局部波长缓变，领先阶给$(S')^2=2m(E-V)=p^2(x)$；下一阶虚部给$2A'p+Ap'=0$。在$p>0$的一个允许区间内，以$x_0$为参考点分离变量，
\[
 \frac{A'}A=-\frac{p'}{2p},\qquad
 \int_{A(x_0)}^{A(x)}\frac{\dd a}{a}
 =-\frac12\int_{p(x_0)}^{p(x)}\frac{\dd q}{q},\qquad
 \ln\frac{A(x)}{A(x_0)}=-\frac12\ln\frac{p(x)}{p(x_0)}.
\]
因此$A=C/\sqrt p$，允许区的两支为
\[
 \varphi_\pm\simeq\frac{C_\pm}{\sqrt{p(x)}}
 \exp\left[\pm\frac i\hbar\int^x p(q)\dd q\right].
\]
检验忽略$A''$的条件，还须把振幅代回原方程。对$A=Cp^{-1/2}$连续求导，
\[
 A'=-\frac C2p^{-3/2}p',\qquad
 A''=\frac{3C}{4}p^{-5/2}(p')^2-\frac C2p^{-3/2}p'',
\]
故相对于领先项$p^2A$的余项为
\[
 \frac{\hbar^2A''}{p^2A}
 =\hbar^2\left[\frac{3(p')^2}{4p^4}-\frac{p''}{2p^3}\right].
\]
一组充分的局部缓变条件是$\hbar|p'|/p^2\ll1$及$\hbar^2|p''|/p^3\ll1$。长区间的累计误差还需结合具体势估计。在转点$p=0$附近条件失效，不能把上述振幅直接延伸穿过转点。

\subsection{平滑转点的四分之一周期相移}
取右侧简单转点$x_t$，即$V(x_t)=E$、$F_t=V'(x_t)>0$，允许区在$x<x_t$。局部线性化$V-E\simeq F_t(x-x_t)$，定态方程成为$\varphi''=(2mF_t/\hbar^2)(x-x_t)\varphi$。令$\ell_t=(\hbar^2/(2mF_t))^{1/3}$、$z=(x-x_t)/\ell_t$，则$\varphi''(x)=\ell_t^{-2}\varphi_{zz}$，代入得$\varphi_{zz}-z\varphi=0$。禁区$z>0$选衰减的$\operatorname{Ai}(z)$。利用艾里函数在$z<0$的渐近连接式，允许区的驻波为
\[
 \varphi(x)\propto\frac1{\sqrt{p(x)}}
 \sin\left[\frac1\hbar\int_x^{x_t}p(q)\dd q+\frac\pi4\right].
\]
这里的相移是局部线性势的连接结果，不能用于突然跳到无限高的硬壁。两个平滑转点各带相同类型修正，闭合相位条件成为$\int_{x_1}^{x_2}p\dd x=(n+1/2)\pi\hbar$，即$J=(n+1/2)h$。

若左端是硬壁$x=0$、右端是平滑转点，将上式连接解代入$\varphi(0)=0$，得到
\[
 \frac1\hbar\int_0^{x_t}p\dd x+\frac\pi4=(n+1)\pi,
\]
所以$J=2\int_0^{x_t}p\dd x=(n+3/4)h$。这给出重力反弹中的$3/4$。两硬壁之间则直接要求站波节点，$\int_0^Lp\dd x=n\pi\hbar$，$J=nh$，从$n=1$开始。三类边界的量子数起点和偏移不同，比较时须把编号一并说明。

\Diagram{fig:m4boundarytypes}{硬壁节点与平滑转点连接所对应的相位条件不同}{\begin{scope}[xshift=-2.4cm]\draw[axis](0,0)--(2.2,0);\draw[very thick](0,-.4)--(0,1.7);\draw[curve,domain=0:2,samples=60]plot(\x,{.9*sin(120*\x)});\node[below]at(1,-.7){硬壁节点};\end{scope}\begin{scope}[xshift=1.2cm]\draw[axis](0,0)--(2.5,0);\draw[curve](0,.7)..controls(.5,1.3)and(.7,1.1)..(1,.5)..controls(1.4,.1)and(1.8,.02)..(2.3,0);\draw[guide](1,-.4)--(1,1.7);\node[below]at(1.2,-.7){平滑转点};\end{scope}}
图\ref{fig:m4boundarytypes}仅区分边界形态。半经典能级公式通常对高量子数较准，因为大部分允许区包含许多波长，转点局部修正占比小。重力反弹可用艾里零点检验误差，谐振子则因二次势的特殊结构使领先WKB能级恰好精确。精确的一个例子不能消除近似条件。

\section{零点能和位置—动量展宽}
对归一化且边界项消失的波函数，定义$\Delta x^2=\langle(x-\langle x\rangle)^2\rangle$、$\Delta p^2=\langle(p-\langle p\rangle)^2\rangle$。取$f=(x-\langle x\rangle)\psi$、$g=(p-\langle p\rangle)\psi$，柯西不等式给$\Delta x^2\Delta p^2\geq|\langle f,g\rangle|^2$。由算符的厄米性及对易关系，
\[
 2i\,\operatorname{Im}\langle f,g\rangle
 =\langle f,g\rangle-\langle g,f\rangle
 =\langle\psi,[x,p]\psi\rangle=i\hbar.
\]
于是$|\langle f,g\rangle|^2\geq(\operatorname{Im}\langle f,g\rangle)^2=\hbar^2/4$，得到
\[
 \Delta x\Delta p\geq\frac\hbar2.
\]
这里的$\Delta x,\Delta p$是同一量子态的概率标准差。位置分布越窄，动量分布的最小展宽越大；方阱基态的非零动能和谐振子零点能都满足这一约束。

对以平衡点为中心、平均动量为零的谐振子态，能量平均满足
\[
 \langle E\rangle=\frac{\Delta p^2}{2m}+\frac{m\omega^2\Delta x^2}2
 \geq\frac{\hbar^2}{8m\Delta x^2}+\frac{m\omega^2\Delta x^2}2.
\]
记$z=\Delta x^2>0$，右侧函数为$f(z)=\hbar^2/(8mz)+m\omega^2z/2$。逐项求导，
\[
 f'(z)=-\frac{\hbar^2}{8mz^2}+\frac{m\omega^2}{2},\qquad
 f''(z)=\frac{\hbar^2}{4mz^3}>0.
\]
令$f'=0$，正根为$z=\hbar/(2m\omega)$，代回得$\langle E\rangle\geq\hbar\omega/2$。可取$\varphi=C\ee^{-b x^2/2}$核实等号态。两次求导为$\varphi'=-bx\varphi$、$\varphi''=(-b+b^2x^2)\varphi$，代入定态方程给
\[
 E\varphi=\left[\frac{\hbar^2b}{2m}
 +\left(\frac{m\omega^2}{2}-\frac{\hbar^2b^2}{2m}\right)x^2\right]\varphi.
\]
定态能量不依赖$x$，故$b=m\omega/\hbar$（取正根以使无穷远衰减），$E=\hbar\omega/2$。零点项同时由量子展宽下界和满足边界条件的波函数确定。

\section{黑体谱中经典近似的边界}
经典均分定理的前提是能量可在$\kb T$尺度上近似连续地取值。对低频模态，能量间隔$h\nu\ll\kb T$，相邻占据数许多，几何分布的离散和可近似为连续积分，于是每模平均$\kb T$。对高频模态，第一次激发就需远大于$\kb T$的能量，绝大多数模态停留在$n=0$，平均能量被玻尔兹曼因子压低。

紫外灾难的根源因此是把低频的连续能量近似用于全部无限多个高频模态。修正只改模态计数是不够的：真空电磁波的态密度$\nu^2$来自空间维数与色散；量子化改变的是各模态的热占据。计算时分别求出态密度和平均热占据，再将二者相乘。

黑体辐射的总能量按$T^4$、光子数按$T^3$变化，平均光子能量因而正比于$T$。把频率谱的横轴换为$h\nu/(\kb T)$、纵轴除以$T^3$，不同温度的曲线重合。实际物体若有选择性吸收、反射、有限模态或非平衡照明，其谱形可偏离这一关系。

\subsection{作用量中的反射相位与量子数编号}
一维束缚运动的作用量为$J(E)=\oint p\dd x=2\int_{x_-}^{x_+}|p(x)|\dd x$。返回支的$p$与$\dd x$同时反号，两支对作用量均给出正贡献。波的闭合相位还包括转点反射相位，因此不同边界不能统一写成$J=nh$。

两个平滑转点时，每个转点引入四分之一周的相位修正，量子条件为$J=(n+1/2)h$，$n=0,1,\ldots$。一硬一软边界给$J=(n+3/4)h$，仍从$n=0$编号；两个无限硬壁则给$J=nh$，应从$n=1$编号。三个式子的$n$从不同数值起算，分别对应各边界问题的节点编号。
\Diagram{fig:m4boundaryphase}{允许区的三类端点；硬壁与平滑转点采用不同反射条件}{\begin{scope}\draw[axis](0,0)--(1.8,0);\draw[curve]plot[smooth]coordinates{(0,1.4)(.3,.5)(.7,.1)(1.1,.5)(1.4,1.4)};\draw[guide](0,.8)--(1.5,.8);\node[below]at(.7,-.2){两软};\end{scope}\begin{scope}[xshift=2.2cm]\draw[axis](0,0)--(1.8,0);\draw[thick](.2,0)--(.2,1.5);\draw[curve](.2,.1)--(1.5,1.4);\draw[guide](.2,.8)--(1.6,.8);\node[below]at(.8,-.2){一硬一软};\end{scope}\begin{scope}[xshift=4.5cm]\draw[axis](0,0)--(1.8,0);\draw[thick](.2,1.5)--(.2,.1)--(1.4,.1)--(1.4,1.5);\draw[guide](.2,.8)--(1.4,.8);\node[below]at(.8,-.2){两硬};\end{scope}}
图\ref{fig:m4boundaryphase}的虚线表示同一类比较中的能量线，交点限定作用量的积分区间。有限高度的陡峭势垒还会产生随能量变化的反射相位，需由波函数与导数的连续条件求出；它通常介于软转点近似与无限硬壁极限之间。

还可把相位条件写成$J(E)=h[n+\eta(E)]$。对量子数求导，
\[
 \left[\frac{\dd J}{\dd E}-h\frac{\dd\eta}{\dd E}\right]\frac{\dd E}{\dd n}=h.
\]
若$\eta$为常数且半经典近似成立，便有$\Delta E\simeq h/T$，其中$T=\dd J/\dd E$是经典周期。若边界相位显著依赖能量，能级间隔还含$\dd\eta/\dd E$修正。无限深方阱可直接解边界方程，重力反弹低能级则应对照艾里函数零点；使用作用量时须区分精确边界解与渐近量子条件。
